\chapter{Cone criteria in arbitrary codimension}\label{sec:rectifiable-cone-characterization}
The reduced-boundary graph argument has a version for an $m$-plane in $\mathbb R^n$, with $1\leq m<n$. The normal component is now a vector in an $(n-m)$-dimensional space. We supply the details needed for the converse implication in \cite{Sim18}, Chapter 3, Theorem 1.6, printed pp.~88--90. A coordinatewise scalar extension is sufficient; an optimal vector-valued extension theorem is not needed.

\begin{definition}[Projection distance and the cone outside a plane]\label{def:plane-projection-cone}
For an $m$-dimensional linear plane $P\subset\mathbb R^n$, let $p_P$ be its orthogonal projection and $q_P=I-p_P$. Put
\[
 d(P,Q)=\|p_P-p_Q\|_{\mathrm{op}},\qquad
 C_\alpha(x,P)=\{z:|q_P(z-x)|\geq\alpha|z-x|\},\quad0<\alpha<1.
\]
The cone consists of directions whose angle from the plane is bounded below. This is equivalent to the complementary-plane convention in \cite{Sim18}, equations (1)--(3), p.~88. For codimension one it agrees with Definition \ref{def:normal-double-cone} after choosing a unit normal.
\end{definition}

\begin{lemma}[Finite nets of planes]\label{lem:grassmann-finite-nets}
The $m$-dimensional planes form a compact metric space under $d$. In particular, for each $\eta>0$ finitely many open $d$-balls of radius $\eta$ cover all of them. Also $\|q_P-q_Q\|_{\mathrm{op}}=d(P,Q)$.
\end{lemma}
\begin{proof}
Identify a plane with its real $n\times n$ orthogonal projection matrix. Such matrices are exactly those satisfying $A^*=A$, $A^2=A$ and $\operatorname{tr}A=m$. These polynomial conditions are closed. A real symmetric idempotent has eigenvalues zero or one, so these matrices are bounded and the trace condition makes their rank exactly $m$. They form a compact subset of the finite-dimensional matrix space. The operator norm induces its usual topology, and compactness gives a finite subcover of the cover by all open balls of radius $\eta$. The equality for complementary projections follows from $q_P-q_Q=-(p_P-p_Q)$. This justifies the finite orientation selection in the cone proof.
\end{proof}

\begin{lemma}[A ball forced outside a nearby plane]\label{lem:plane-cone-ball-inclusion}
Let $x\ne y$, $d_0=|y-x|$, and let $P,P_x$ be $m$-planes with $d(P,P_x)\leq1/8$. If $|q_P(y-x)|>d_0/2$, then
\[
 B_{d_0/32}(y)\subset C_{1/4}(x,P_x)\cap B_{2d_0}(x).
\]
This is the vector-normal version of Lemma \ref{lem:normal-cone-ball-inclusion}.
\end{lemma}
\begin{proof}
The projection estimate gives $|q_{P_x}(y-x)|>3d_0/8$. For $|z-y|<d_0/32$, use $\|q_{P_x}\|_{\mathrm{op}}=1$ to obtain $|q_{P_x}(z-x)|>11d_0/32$, while $|z-x|<33d_0/32<2d_0$. Their ratio exceeds $1/3$ and hence $1/4$. No sign choice for a normal vector is involved, so the calculation works in every codimension.
\end{proof}

\begin{proposition}[Uniform plane-cone bounds imply a vector-valued graph]\label{prop:uniform-plane-cone-graph}
Let $\mu$ be Radon on an open $U\subset\mathbb R^n$ and $S\subset U$. Suppose $c,\rho>0$, an $m$-plane $P$, and $m$-planes $P_x$ for $x\in S$ satisfy $d(P_x,P)\leq1/8$, $B_\rho(x)\subset U$, and, for every $0<r<\rho$,
\[
 \mu(B_r(x))\geq cr^m,\qquad
 \mu(C_{1/4}(x,P_x)\cap B_r(x))\leq\delta r^m,
 \qquad \delta=\frac{c}{2\cdot64^m}.
\]
If $\operatorname{diam}S<\rho/2$, then $S$ lies on the graph over $P$ of a map $h:P\to P^\perp$ with Lipschitz constant at most $\sqrt{(n-m)/3}$. On the projected subset $p_P(S)$ the original height map has the sharper bound $1/\sqrt3$.
\end{proposition}
\begin{proof}
If a pair $x,y\in S$ had $|q_P(y-x)|>|y-x|/2$, Lemma \ref{lem:plane-cone-ball-inclusion} would give
\[
 c(d_0/32)^m\leq\mu(B_{d_0/32}(y))
 \leq\mu(C_{1/4}(x,P_x)\cap B_{2d_0}(x))
 \leq\delta(2d_0)^m=\tfrac12c(d_0/32)^m,
\]
where $d_0=|y-x|<\rho/2$. This is impossible. Hence every pair satisfies $|q_P(y-x)|\leq|y-x|/2$. Orthogonality gives $|p_P(y-x)|\geq\sqrt3|y-x|/2$ and $|q_P(y-x)|\leq|p_P(y-x)|/\sqrt3$. Projection is injective, and the height map on its image is $h_0(p_Px)=q_Px$ with the claimed bound.

Choose an orthonormal basis of $P^\perp$. Each of the $n-m$ components of $h_0$ is scalar $1/\sqrt3$-Lipschitz. Proposition \ref{prop:mcshane-extension} extends it to all of $P$ with that bound. The Euclidean norm of the vector of component differences is at most $\sqrt{(n-m)/3}$ times the base distance. These extensions form $h$. If $S$ is empty take $h=0$. The resulting whole graph is a Lipschitz image of $\mathbb R^m$; the possibly enlarged extension constant has no effect on containment or rectifiability.
\end{proof}

\begin{theorem}[Pointwise plane-cone bounds give a countable graph cover]\label{thm:pointwise-plane-cone-cover}
Let $\mu$ be Radon on an open $U\subset\mathbb R^n$, and let $S\subset U$. Suppose that for each $x\in S$ there is an $m$-plane $P_x$ for which
\[
 \liminf_{r\downarrow0}\frac{\mu(B_r(x))}{r^m}>0,\qquad
 \lim_{r\downarrow0}\frac{\mu(C_{1/4}(x,P_x)\cap B_r(x))}{r^m}=0.
\]
Then $S$ is contained in countably many whole Lipschitz graphs over $m$-planes. No measurability of the plane selection is needed for this containment assertion.
\end{theorem}
\begin{proof}
For positive integers $k,j$, retain those points for which $B_{1/j}(x)\subset U$ and the lower bound with $c=1/k$ and cone upper bound with $\delta=c/(2\cdot64^m)$ both hold for all $0<r<1/j$. The two limits imply that these countably many classes cover $S$: first choose a smaller positive lower constant, then a sufficiently small common scale at the given point. Split each class using the finite $1/8$-net of planes from Lemma \ref{lem:grassmann-finite-nets}, allowing overlap. Intersect with a countable cover of space by balls of diameter less than $1/(2j)$. Each resulting piece satisfies Proposition \ref{prop:uniform-plane-cone-graph}. There are countably many pieces and hence countably many graphs. This construction uses uniform inequalities only on the defined pieces, not an unsupported uniform convergence assertion on $S$ itself.
\end{proof}

\begin{lemma}[A plane tangent supplies density and cone decay]\label{lem:plane-tangent-density-cone}
Let $\mu$ be Radon on an open $U\subset\mathbb R^n$, let $x\in U$, and suppose
\[
 r^{-m}(\eta_{x,r})_\#\mu\rightharpoonup\mathcal H^m\llcorner P,
 \qquad \eta_{x,r}(y)=(y-x)/r,
\]
locally as $r\downarrow0$, where $m\geq1$. Then $\mu(B_r(x))/(\omega_mr^m)\to1$ and, for every $0<\alpha<1$,
\[
 r^{-m}\mu(C_\alpha(x,P)\cap B_r(x))\to0.
\]
This is the content of \cite{Sim18}, Chapter 3, Remark 1.5(3), p.~87.
\end{lemma}
\begin{proof}
The ball mass follows by continuous cutoffs as in Corollary \ref{cor:c1-patch-unit-density}. For the cone, its rescaling inside $B_1$ lies in the compact set $K=\{|q_Pz|\geq\alpha|z|\}\cap\overline B_1$. Its intersection with $P$ is only the origin, which has zero $\mathcal H^m$ measure for $m\geq1$. The upper bound for compact sets under weak Radon convergence gives a limiting upper mass of zero. These two arguments use different continuity facts: the unit sphere has zero plane measure, whereas the closed cone compact set itself has zero plane measure.
\end{proof}

\begin{theorem}[Almost-everywhere plane tangents imply rectifiability]\label{thm:rectifiable-converse-tangents}
Let $M\subset\mathbb R^n$ be $\mathcal H^m$-measurable and locally finite in $\mathcal H^m$ measure, with $1\leq m<n$. If $T_xM$ in Definition \ref{def:approximate-tangent} exists at almost every point of $M$, then $M$ is countably $m$-rectifiable. This proves the converse in \cite{Sim18}, Chapter 3, Theorem 1.6, pp.~88--90.
\end{theorem}
\begin{proof}
Use Lemma \ref{lem:rectifiable-borel-representative} to regard $\mu=\mathcal H^m\llcorner M$ as a locally finite Radon measure. Let $S\subset M$ be the full-measure set on which the tangent exists, and assign $P_x=T_xM$. Lemma \ref{lem:plane-tangent-density-cone} supplies both hypotheses of Theorem \ref{thm:pointwise-plane-cone-cover}, with positive density $\omega_m$. Thus $S$ lies in countably many Lipschitz graphs, each a Lipschitz image of $\mathbb R^m$. The remaining subset of $M$ has zero $\mathcal H^m$ measure, exactly the allowed exception in Definition \ref{def:rectifiable-set}. Nothing in this argument identifies a merely positive finite density with a plane tangent; the directional cone decay is also required.
\end{proof}
